\chapter{Objetivo}

El objetivo de este proyecto es el diseño e implementación de una plataforma experimental de bajo coste orientada a la emulación de instrumentación aviónica. Para ello el sistema está basado en un microcontrolador ESP32 como núcleo, y varios sensores, entre los cuales se encuentran sensores inerciales, barómetro, telemetría, radar, entre otros, a decidir en función del desarrollo del proyecto. Con la información proporcionada por los sensores se pretende calcular la actitud, el rumbo, la altitud, la velocidad vertical de la plataforma en tiempo real. Por último también se desarrollará el software para visualizar estos datos en una pantalla similar a la de una PFD (Primary Flight Display). \newline
 
El resultado final deseado es un sistema integrado por 3 subsistemas principales: (1) una PCB que integre los componentes físicos del proyecto, como el controlador y los sensores. (2) el firmware que recibe los datos de los sensores, los integra con fusión de datos y envía los resultados al tercer subsistema. (3) Un PFD implementado con el lenguaje Processing que muestre los valores calculados. 



